\section{Successor interface: failure episodes, scoped failure knowledge, and atomized invention gates}
\label{sec:p5-successor-failure-discovery}

The successor ORION failure programme separates two records that are easy to conflate. A \emph{failure episode} preserves what happened: source or execution identity, chronology, observed failure mode, propagated effect, candidate causes, recovery attempts, residuals and the evidence available at the time. \emph{Failure knowledge} is a stronger derived object: it states which conclusion may be reused, the evidence and responsibility attribution supporting that conclusion, the context coordinates that must remain the same, and the changes that reopen the question. The episode can exist without the stronger reusable exclusion.

This distinction makes the recurrence--cause boundary operational. Repeatedly observing the same signature can justify a recurrence hypothesis, but a reusable method lesson still needs responsibility evidence and an applicability contract. Every load-bearing context coordinate must have an explicit role. A required-same mismatch makes the old exclusion inapplicable; a declared semantic or representation change may reopen it; an unclassified context cannot make a failure globally applicable by omission. Reopening removes an old prohibition only for the changed regime---it is not positive evidence that the new method works.

The merged prospective failure-memory study preserves these semantics but narrows the architectural claim. Its earned terminal is \texttt{FAILURE\_MEMORY\_USEFUL\_BUT\_STANDARD\_METHODS\_SUFFICIENT}: scoped negative history is useful, yet the tested benefit is absorbed by strong standard failure-learning/memory parents once context, responsibility and reopen conditions are represented. Self-ORION therefore retains the scoped record and removes any unsupported claim that a distinct generic ORION failure-memory mechanism is required.

Before Self-ORION treats a residual as an invention opportunity, the successor programme also decomposes broad labels such as epistemic tension, imagination and concept formation into bounded candidate atoms, mature parent mechanisms and registered interactions. The three pilots close as \texttt{MULTIPLE\_TENSION\_ATOMS\_REQUIRED}, \texttt{MULTIPLE\_IMAGINATION\_ATOMS\_REQUIRED} and \texttt{MULTIPLE\_CONCEPT\_ATOMS\_REQUIRED}; the shared outcome is a frozen recursive atom-study calculus rather than three new universal modules. Historical morphology and failure atlases organize challenger families and expose discriminators, but remain design evidence rather than protected gold.

The protected zero-error Jump study supplies the strongest adoption lesson. ORION-JUMP and the verified representation-regime revision parent tie exactly on both disjoint frozen splits, so the registered terminal is \texttt{REPRESENTATION\_INVENTION\_NO\_INCREMENTAL\_VALUE}. Under this paper's governance, the correct response to that result is parent subsumption: retain the verified parent and the bounded testing interface, and do not manufacture a distinct Jump operator merely to preserve an ORION label.

The adoption path therefore remains the one defined in this paper: diagnose the issue, search and structurally absorb known mechanisms, construct a challenger, execute it in isolation, require replay and fresh-transfer evidence, pass protected assurance, and leave final disposition to the external host. None of the successor results supplies H1--H4 evidence for governed self-improvement; those hypotheses remain \texttt{CANNOT\_CHECK}.